\chapter{Planteamiento del problema, hipótesis y justificación}\label{chap:planteamiento}
\begin{sloppypar}

\section{Planteamiento del problema}

Dado el modelo físico del sistema de levitación magnética ECP-730 en configuración atractiva, inestable en lazo abierto y con efecto atascamiento-deslizamiento, se busca diseñar un controlador de posición con las herramientas del control lineal clásico. Se propone un controlador con doble acción integral (\textbf{PII}), diseñado sobre el modelo linealizado y evaluado sobre el modelo no lineal con fricción estática y zona muerta.

\section{Hipótesis}

Un controlador lineal con doble acción integral, diseñado sobre el modelo linealizado del sistema de levitación magnética en configuración atractiva:
\begin{enumerate}
  \item estabiliza la posición del levitante en presencia de la fricción estática y de la zona muerta del actuador, sin necesidad de compensarlas explícitamente, y
  \item por la alta ganancia que proporciona a baja frecuencia, reduce el tiempo de atascamiento y la amplitud de las oscilaciones producidas por el atascamiento-deslizamiento respecto a un controlador con una sola acción integral.
\end{enumerate}
La segunda parte retoma la observación de Pérez-Gómez \cite{perez} en un motor de CD y la de Hernández Alcántara \cite{diana_tesis}, según la cual el sistema de control debe proporcionar alta ganancia en estado estacionario para reducir los efectos de la fricción y la zona muerta.

\section{Justificación}

Los sistemas de levitación magnética con componentes mecánicos en contacto son sistemas de ingeniería complejos que han interesado a la comunidad científica durante años. La zona muerta, una de las no linealidades no suaves, es un reto para los controladores lineales: puede causar inestabilidad y deteriora el desempeño en estado estacionario en lazo cerrado, sobre todo alrededor del punto de equilibrio, donde los desplazamientos son pequeños y el sistema opera dentro de ella. Allí provoca oscilaciones de alta frecuencia en la variable de control, y si a esto se añade la dinámica no lineal de los electroimanes, el diseño del esquema de control se vuelve difícil.

Trabajos como \cite{khimani,mukesh,neelma} no incluyen la zona muerta en su modelado, ya sea porque la ignoran o porque la tratan como dinámica no modelada o como incertidumbre no estructurada. Las estrategias de control propuestas para este problema van desde controladores robustos y adaptativos difusos y QFT hasta controladores PID \cite{kuo1,mukesh,maaz}, que manejan bien el error en estado estacionario, pero la literatura sigue siendo escasa. En cuanto a la representación de la zona muerta, se han propuesto modelos lineales \cite{salim}, no lineales \cite{huanqing}, difusos \cite{zhi}, asimétricos \cite{gang}, simétricos \cite{nizar} y generalizados \cite{guanyu}, entre otros; sin embargo, estos modelos y los resultados experimentales que los acompañan no representan por completo el comportamiento real.

Un controlador lineal de estructura sencilla, diseñado con técnicas clásicas, es atractivo por su facilidad de análisis e implementación. Evaluar si un controlador de este tipo estabiliza el sistema de levitación en su configuración inestable y qué desempeño alcanza frente a la fricción y a la zona muerta dinámica, que depende de la posición del cuerpo levitante, permite establecer el alcance de las estrategias lineales en este problema y orientar el diseño de estrategias más elaboradas.

\end{sloppypar}
